\section{Methodology}
\label{sec:method}

\subsection{Publication Model and Correctness Criteria}
\label{sec:model}

\paragraph{Producer and monitor.}
At control step $k$ the producer publishes a 64-byte witness record
\[
w_k = (\mathit{epoch},\ \mathit{ts},\ d_a,\ d_b,\ d_c,\ \mathit{meta}),
\]
where $d_a,d_b,d_c$ are three-phase actuation values and $\mathit{meta}$
carries a configuration identifier, a reserved MAC-sized payload, and a
redundant low-16-bit copy of $\mathit{epoch}$. The monitor attempts to
read a snapshot $\hat w_k$, aligns it to an independently acquired
evidence sample $e_k = (\hat d_a, \hat d_b, \hat d_c)$, and computes the
residual
\[
r_k \;=\; \max_i \bigl| d_i(\hat w_k) - \hat d_i(e_k) \bigr|,
\]
raising an alarm when $r_k>\varepsilon$. A small bounded-window \STL
policy is evaluated over the residual stream so that the consumer
exercises the signal path realistically.

\paragraph{Torn records as a correctness failure.}
The torn-read detector compares the low 16 bits of $\mathit{epoch}$
against the redundant copy stored inside $\mathit{meta}$. Any mismatch is
a \emph{torn read}: the monitor has observed at least two different
publications interleaved into one logical read. This is not a minor
performance artifact. It is a correctness failure, because the policy
layer is now reasoning about a state the system was never in.

\paragraph{Publication contract.}
The publication channel should satisfy the following contract:
\begin{quote}
\emph{Every successful monitor read corresponds to exactly one producer
publication. If the reader cannot establish that property, it must return
a retry signal rather than data.}
\end{quote}
Coherence is part of the substrate on which this contract is built, but
it does not satisfy the contract by itself. The five architectures in
\S\ref{sec:protocols} are five concrete ways to either satisfy or
deliberately violate it.

\paragraph{Comparison logic.}
The matrix is constructed so that one factor changes at a time.
Comparing \unsync, \seqlock, and \dblbuf holds transport fixed and asks
what the publication primitive alone contributes on a coherent shared
record. Comparing \dmanaive and \dmaseqlock holds the explicit-transfer
path fixed and asks the same question on the mirror side. Comparing
\seqlock with \dmaseqlock then isolates direct coherent publication
against explicit transfer under the same odd-even validity rule.
\dblbuf adds a stronger direct-sharing primitive designed for repeated
publication rather than exclusion alone. This structure lets us
attribute differences to the primitive, the transport, or their
interaction rather than to unrelated benchmark changes.

\begin{figure}[t]
  \centering
  \includegraphics[width=\columnwidth]{figures/witness_evidence_conformance_arch.pdf}
  \caption{System model for the study. The producer publishes witness
  records into the evaluated channel, the monitor reads snapshots from
  that channel and checks them against an independently acquired
  evidence stream, and only the producer-to-monitor channel is treated
  as the architectural object of study.}
  \label{fig:arch}
\end{figure}
